% Chapter 2 -- The Calculator (feature topic).  Sources (Law 1):
% os/apps/calcda/calcda.c (CALCULAT.EXE: the rounded window titled
% "Calculator", the File menu Close ^W, the Edit menu all dimmed; a click
% on a key and the key map end in the one calc_press); calc_core.h and
% calc_core.c (the engine, from 101 probes of the real Mac OS 8.1
% Calculator: chain arithmetic without precedence, an operator replaces an
% operator, '=' does not repeat, 13 characters of entry, the 12-character
% readout, infinity / -infinity / not a number, C / c / NumLock clear,
% Enter is '=', Escape / Backspace do nothing; calc_cmd_from_key);
% spec/calcda/calcda_layout.h (18 keys, the readout).  The menu entries
% are registered against CALCULAT's own bar (flair_app_menus.c).
% GATES: test-calcda test-flair-calcda test-flair-apple-menu test-pubs-flair test-pubs-flair-screens
\Gates{test-calcda, test-flair-calcda, test-flair-apple-menu, test-pubs-flair,
  test-pubs-flair-screens}
\chapter{The Calculator}
\label{ch:calculator}
\index{Calculator}

The Calculator works like a pocket calculator.  You can use it while you
work in another window, and copy its arithmetic onto paper or into a
document.

\section{Opening the Calculator}
\Operation{APPLE-LAUNCH}
\begin{TaskSteps}
\Step Choose \menuitem{Calculator} from the Apple menu.
\Result{The Calculator opens in a small window with rounded corners, near
the top of the screen.  The lower menu bar changes to the Calculator's
menus.}
\end{TaskSteps}

\FlairScreen[scale=4, label=fig:calc]{a_calc}{The Calculator}

You can also open it by double-clicking \file{CALCULAT.EXE} in the
\file{APPS} folder.  Like every application from the disk, it does not
open while another one is open.

\section{The parts of the Calculator}
\SideHead{Display}The box at the top shows the number you are typing, or
the result.  It holds 12 characters.
\SideHead{Keys}Eighteen keys: the digits \texttt{0} to \texttt{9}, the
decimal point, the four operations \texttt{+ - * /}, \texttt{=}, and
\texttt{C}, which clears the Calculator.
\SideHead{Close box}Closes the Calculator.  It does not ask anything: the
Calculator keeps nothing that needs saving.

\section{Working out a sum}
\index{Calculator!using}
\begin{TaskSteps}
\Step Type \type{12+30=} on the keyboard.
\Result{The display shows \texttt{42}.}
\end{TaskSteps}

\FlairScreen[scale=4, label=fig:calcsum]{a_calcsum}{12 + 30 = 42}

Clicking the Calculator's keys does exactly what typing does.

\begin{TaskSteps}[resume]
\Step Click \texttt{7}, \texttt{+}, \texttt{8} and \texttt{=} on the
Calculator.
\Result{The display shows \texttt{15}.}
\end{TaskSteps}

\FlairScreen[scale=4, label=fig:calcclick]{a_calcclick}{7 + 8 = 15, clicked}

\section{The keys on the keyboard}
\index{Calculator!keys}
\begin{center}\small
\begin{tabular}{@{}p{38mm}p{61mm}@{}}
\textbf{Press} & \textbf{To}\\\hline
\key{0} to \key{9}, \key{.} & enter a number\\
\key{+} \key{-} \key{*} \key{/} & add, subtract, multiply, divide\\
\key{=} or \key{Enter} & see the result\\
\key{C} or \key{Num Lock} & clear the Calculator\\
\end{tabular}
\end{center}
Other keys, \key{Esc} and \key{Backspace} among them, do nothing.  To
correct a number you have typed, clear the Calculator and start the
calculation again.

\section{How the Calculator works out results}
\index{Calculator!order of operations}
The Calculator carries out each operation as you go, from left to right.
It does not do multiplication and division before addition and
subtraction: typing \type{2+3*4=} shows \texttt{20}, because 2 + 3 is
worked out first and then multiplied by 4.

\FlairScreen[scale=4, label=fig:calcchain]{a_calcchain}{2 + 3 * 4 = 20}

\SideHead{Two operations}If you type one operation right after another,
the second replaces the first: \type{5+*2=} shows \texttt{10}.
\SideHead{The equals key}\key{=} finishes the calculation.  Pressing it
again does not repeat the last operation.  Pressing it straight after an
operation drops that operation: \type{2+=} shows \texttt{2}.
\SideHead{After a result}Typing a digit or a point after \key{=} starts a
new number; typing an operation goes on from the result.
\SideHead{Long numbers}A number you type can have up to 13 characters,
the point and a leading \texttt{0.} included; more digits are ignored.
Results that do not fit the 12-character display are shown in powers of
ten, such as \texttt{1E12} for 1,000,000,000,000 and \texttt{3.3333333E-8}
for 0.000000033333333.
\SideHead{Dividing by zero}\type{1/0=} shows \texttt{infinity},
\type{0-1/0=} shows \texttt{-infinity}, and \type{0/0=} shows \texttt{not a
number}.  Clear the Calculator, or type a digit, to go on.

\FlairScreen[scale=4, label=fig:calcinf]{a_calcinf}{1 / 0 =}

\section{Everyday calculations}
\index{Calculator!examples}
\SideHead{Adding a column}Type each amount followed by \key{+}, and finish
the last with \key{=}.  To add 12.50, 7.25 and 30, type
\type{12.5+7.25+30=}.  The display shows \texttt{49.75}.

\FlairScreen[scale=4, label=fig:calccol]{a_calccol}{12.5 + 7.25 + 30 =}

\SideHead{A share as a percentage}Divide the part by the whole, then
multiply by 100.  To find what share 450 is of 1200, type
\type{450/1200*100=}.  Because the Calculator works from left to right, the
division is done first, as you want, and the display shows \texttt{37.5}.

\FlairScreen[scale=4, label=fig:calcpct]{a_calcpct}{450 / 1200 * 100 =}

\SideHead{Starting again}Press \key{C}, or click \texttt{C}, before each
new calculation.  A digit typed straight after a result also starts a new
number.

\section{The Calculator's menus}
\FlairScreen[scale=4, label=fig:calcfile]{a_calcfile}{The Calculator's File
menu}
\subsection{File}
\begin{MenuReference}{CALCULAT}{File}
\MenuItemEntry{Close}{W}{on}{Closes the Calculator.}
\end{MenuReference}
\FlairScreen[scale=4, label=fig:calcedit]{a_calcedit}{The Calculator's Edit
menu}
\subsection{Edit}
\begin{MenuReference}{CALCULAT}{Edit}
\MenuItemEntry{Undo}{Z}{off}{Dimmed.}
\MenuItemEntry{Cut}{X}{off}{Dimmed.}
\MenuItemEntry{Copy}{}{off}{Dimmed.  The Calculator does not put its
display on the Clipboard in this version; write the result down, or type
it into your document.}
\MenuItemEntry{Paste}{V}{off}{Dimmed.}
\MenuItemEntry{Clear}{}{off}{Dimmed.  Use the \texttt{C} key.}
\end{MenuReference}
